% Complete Section 4.4, original PDF 49-54 / printed 39-44.
% The user-approved non-Fraktur mappings preserve harmonic/cohomology distinctions.
% block: abstract-hodge-heading page:49 type:section
\subsection{{The abstract Hodge Laplacian}}\label{sec:abstract-hodge-laplacian}

% block: hodge-laplacian-domain-intro page:49 type:ordinary-paragraph
{Let $(W,d)$ be a closed Hilbert complex. Associated to it is a graded operator from $W\to W$ of degree $0$ given by $L=dd^*+d^*d$. In the case of the de Rham complex associated to a Riemannian manifold, this operator is called the Hodge Laplacian, and we refer to it generally as the (abstract) Hodge Laplacian associated to the complex. Specifically,}
\begin{equation}\label{eq:abstract-hodge-domain}
D(L^k)=\{\,u\in V^k\cap V_k^*\mid du\in V_{k+1}^*,\ d^*u\in V^{k-1}\,\},
\end{equation}
{and}
\begin{equation}\label{eq:abstract-hodge-operator}
\begin{gathered}
L^k u=d_{k+1}^*d^k u+d^{k-1}d_k^*u=(d^k)^*d^k u+d^{k-1}(d^{k-1})^*u,\qquad u\in D(L^k).
\end{gathered}
\end{equation}

% block: hodge-kernel-identification page:49 type:ordinary-paragraph
{A major subject of FEEC is the numerical solution of the Hodge Laplace equation $L^k u=f$ for given $f\in W^k$. Whether this problem has a unique solution depends on whether the space $\mathsf H^k$ of harmonic forms vanishes. Indeed, we now show that the null space of $L^k$ is precisely $\mathsf H^k$, explaining the terminology \emph{harmonic}. From \eqref{eq:harmonic-cocycle-orthogonality} we see that $\mathsf H^k\subset D(L^k)$ and $L^k u=0$ for all $u\in\mathsf H^k$. Thus $\mathsf H^k\subset\mathcal N(L^k)$. On the other hand, from the definition of $L$ we have that}
\begin{equation}\label{eq:hodge-green-identity}
\langle L^k u,v\rangle=\langle du,dv\rangle+\langle d^*u,d^*v\rangle,\qquad u\in D(L^k),\ v\in V^k\cap V_k^*.
\end{equation}
{Taking $v=u$, we get}
\[
\|du\|^2+\|d^*u\|^2=\langle L^k u,u\rangle,\qquad u\in D(L^k),
\]
{and so,}
\[
\|du\|^2+\|d^*u\|^2=0,
\]
{if $u\in\mathcal N(L^k)$. Thus, if $u\in\mathcal N(L^k)$, then $du$ and $d^*u$ vanish, and so $u\in\mathsf H^k$. This verifies that $\mathcal N(L^k)=\mathsf H^k$, i.e., that solutions of the Hodge Laplacian are determined only up to addition of a harmonic form.}

% block: hodge-solvability-orthogonality page:49 type:ordinary-paragraph
{The harmonic forms also determine a necessary condition for the existence of solutions. Taking the inner product of the equation $L^k u=f$ with $q\in\mathsf H^k$ and using \eqref{eq:hodge-green-identity}, we see that $\langle f,q\rangle=0$. Thus a necessary condition for the existence of a solution is the orthogonality condition $f\perp\mathsf H^k$. In fact, in Theorem~\ref{thm:hodge-well-posedness} below we shall show that this condition is not only necessary but also sufficient for existence of a solution.}

% block: hodge-three-formulations-heading page:50 type:subsection
\subsubsection{{Three formulations of the Hodge Laplace problem}}\label{subsec:hodge-three-formulations}

% block: hodge-strong-form-intro page:50 type:ordinary-paragraph
{In order to formulate a problem which is solvable for any $f\in W^k$, we make use of the orthogonal projection $P_{\mathsf H}$ of $W^k$ onto $\mathsf H^k$. Then we can formulate the following \emph{strong form} of the Hodge Laplace equation, which we shall show is well-posed.}

% block: hodge-strong-formulation page:50 type:ordinary-paragraph
\medskip\noindent\textbf{{Strong formulation}}\par
{Given $f\in W^k$, find $u\in D(L^k)$ such that}
\begin{equation}\label{eq:hodge-strong-formulation}
Lu=f-P_{\mathsf H}f,\qquad u\perp\mathsf H^k.
\end{equation}

% block: hodge-solution-decomposition page:50 type:ordinary-paragraph
{Note that if $u$ is a solution to this problem, then}
\begin{equation}\label{eq:hodge-data-decomposition}
f=dd^*u+P_{\mathsf H}f+d^*du.
\end{equation}
{Clearly $dd^*u\in\mathcal B$, $P_{\mathsf H}f\in\mathsf H$, and $d^*du\in\mathcal B^*$, so \eqref{eq:hodge-data-decomposition} is the Hodge decomposition of $f$, i.e., $dd^*u=P_{\mathcal B}f$ and $d^*du=P_{\mathcal B^*}f$.}

% block: hodge-weak-formulations-intro page:50 type:ordinary-paragraph
{We will also consider two weak formulations of this problem. The primal formulation corresponds directly to the strong form.}

% block: hodge-primal-formulation page:50 type:ordinary-paragraph
\medskip\noindent\textbf{{Primal weak formulation}}\par
{Given $f\in W^k$, find $u\in V^k\cap V_k^*$ such that $u\perp\mathsf H^k$ and}
\begin{equation}\label{eq:hodge-primal-formulation}
\langle du,dv\rangle+\langle d^*u,d^*v\rangle=\langle f-P_{\mathsf H}f,v\rangle,\qquad v\in V^k\cap V_k^*.
\end{equation}

% block: hodge-mixed-variables page:50 type:ordinary-paragraph
{For the mixed weak formulation, we introduce two new variables, $\sigma=d^*u\in V^{k-1}$ and $p=P_{\mathsf H}f\in\mathsf H^k$.}

% block: hodge-mixed-formulation page:50 type:ordinary-paragraph
\medskip\noindent\textbf{{Mixed weak formulation}}\par
{Given $f\in W^k$, find $\sigma\in V^{k-1}$, $u\in V^k$, and $p\in\mathsf H^k$ such that}
\begin{equation}\label{eq:hodge-mixed-formulation}
\begin{aligned}
\langle\sigma,\tau\rangle-\langle u,d\tau\rangle&=0,&\tau&\in V^{k-1},\\
\langle d\sigma,v\rangle+\langle du,dv\rangle+\langle p,v\rangle&=\langle f,v\rangle,&v&\in V^k,\\
\langle u,q\rangle&=0,&q&\in\mathsf H^k.
\end{aligned}
\end{equation}

% block: hodge-mixed-asymmetry page:50 type:ordinary-paragraph
{An important observation is that the mixed weak formulation, unlike the other two formulations, does not treat the operators $d$ and $d^*$ symmetrically. Only the $d$ operator appears in this formulation, the terms involving $d^*$ being treated weakly.}

% block: hodge-formulation-equivalence-intro page:50 type:ordinary-paragraph
{Despite their different appearance, the three formulations are completely equivalent.}

% block: hodge-formulation-equivalence page:50 type:theorem
\Needspace{8\baselineskip}
\begin{Theorem}{{equivalence of formulations}}\label{thm:hodge-formulation-equivalence}
{Let $f\in W^k$. An element $u\in W^k$ solves the strong formulation of the abstract Hodge Laplace problem if and only if it solves the primal weak formulation. Moreover, in this case, if we set $\sigma=d^*u$ and $p=P_{\mathsf H}f$, then the triple $(\sigma,u,p)$ solves the mixed weak formulation. Finally, if some $(\sigma,u,p)$ solves the mixed weak formulation, then $u$ solves the other two formulations of the problem, and $\sigma=d^*u$, $p=P_{\mathsf H}f$.}
\end{Theorem}

% block: hodge-equivalence-proof-regularity page:51 type:proof
\begin{Proof}
{First we show the equivalence of the strong formulation and the primal weak formulation. If $u$ satisfies the strong formulation, then certainly $u\in V\cap V^*$ and satisfies the primal weak formulation. On the other hand, if $u$ satisfies the primal weak formulation and \emph{if $u$ belongs to $D(L)$}, then it satisfies the strong formulation. Therefore, we need to show that any $u$ which satisfies the primal weak formulation, and so a priori only belongs to $V\cap V^*$, actually belongs to $D(L)$ defined in \eqref{eq:abstract-hodge-domain}, i.e., satisfies $du\in V^*$ and $d^*u\in V$. The statement that $du\in V^*$ means that there exists a constant $c$ such that}
\begin{equation}\label{eq:hodge-du-adjoint-bound}
\langle du,dv\rangle\leq c\|v\|,\qquad v\in V.
\end{equation}

% block: hodge-equivalence-proof-hodge-split page:51 type:proof
{Using the Hodge decomposition, we split $v=v_1+v_2$ with $v_1\in\mathcal Z$, $v_2\in V\cap\mathcal B^*$. Obviously,}
\[
\langle du,dv_1\rangle=0,
\]
{while}
\[
\langle du,dv_2\rangle=\langle du,dv_2\rangle+\langle d^*u,d^*v_2\rangle=\langle f-P_{\mathsf H}f,v_2\rangle\leq c\|v_2\|\leq c\|v\|,
\]
{where $c=\|f-P_{\mathsf H}f\|$. Adding these results gives \eqref{eq:hodge-du-adjoint-bound} and shows that $du\in V^*$. An analogous argument shows that $d^*u\in V$, and so we have proven that $u\in D(L)$ and so that a solution of the primal weak formulation is a strong solution.}

% block: hodge-equivalence-proof-strong-mixed page:51 type:proof
{Next, suppose $u$ solves the strong formulation, and set $\sigma=d^*u$, $p=P_{\mathsf H}f$. Since $u\in D(L)$, $\sigma\in V$. The first equation of the mixed formulation is a consequence of the definition of $\sigma$. We have, from the Hodge Laplace equation, that}
\[
d\sigma+d^*du+p=f,
\]
{and, taking the inner product with any $v\in V$, we obtain the second equation of the mixed method. The third equation follows from the assumption that $u\perp\mathsf H^k$. This shows that $(\sigma,u,p)$ solves the weak formulation.}

% block: hodge-equivalence-proof-mixed-primal page:51 type:proof
{Finally, if $(\sigma,u,p)\in V\times V\times\mathsf H$ solves the mixed weak formulation, then the first equation implies $u\in V^*$ and $d^*u=\sigma$, and the last equation implies that $u\perp\mathsf H$. Taking $v\in\mathsf H$ in the second equation shows that $p=P_{\mathsf H}f$. Finally, taking $v\in V\cap V^*$ in the second equation we conclude that $u$ solves the primal weak formulation.}
\end{Proof}

% block: hodge-well-posedness-heading page:51 type:subsection
\subsubsection{{Well-posedness}}\label{subsec:hodge-well-posedness}

% block: hodge-well-posedness-intro page:51 type:ordinary-paragraph
{Finally, we show that the abstract Hodge Laplace equation is well-posed once we account for any harmonic forms.}

% block: hodge-well-posedness page:51 type:theorem
\Needspace{8\baselineskip}
\begin{Theorem}{{well-posedness}}\label{thm:hodge-well-posedness}
{For each $f\in W^k$ there exists a unique solution $u\in D(L^k)$ such that $u\perp\mathsf H^k$ and $Lu=f-P_{\mathsf H}f$. Moreover}
\[
\|u\|+\|du\|+\|d^*u\|+\|dd^*u\|+\|d^*du\|+\|p\|\leq c\|f\|.
\]
{The constant $c$ entering this estimate depends only on the constant $c_P$ entering the Poincaré inequality (Theorem~\ref{thm:hilbert-complex-poincare}).}
\end{Theorem}

% block: hodge-well-posedness-proof-reduction page:52 type:proof
\begin{Proof}
{When a solution exists, the bounds on $\|dd^*u\|$, $\|d^*du\|$, and $\|p\|$ are immediate from the Hodge decomposition $f=dd^*u+d^*du+p$, so it suffices to show that for any $f\in W$ a unique solution exists and satisfies}
\[
\|u\|_V+\|d^*u\|\leq c\|f\|.
\]

% block: mixed-hodge-bilinear-form page:52 type:proof
{Our proof is based on the mixed weak formulation. Writing $X=V^{k-1}\times V^k\times\mathsf H^k$ for brevity, define a bounded bilinear form $B:X\times X\to\mathbb R$ by combining the left-hand sides of the three equations defining the mixed formulation:}
\begin{equation}\label{eq:hodge-mixed-bilinear-form}
\begin{aligned}
B(\sigma,u,p;\tau,v,q)={}&\langle\sigma,\tau\rangle-\langle u,d\tau\rangle-\langle d\sigma,v\rangle\\
&-\langle du,dv\rangle-\langle p,v\rangle-\langle u,q\rangle.
\end{aligned}
\end{equation}

% block: mixed-hodge-load-functional page:52 type:proof
{Note that $B$ is symmetric. Also define a bounded linear functional $F:X\to\mathbb R$ by}
\begin{equation}\label{eq:hodge-mixed-load-functional}
F(\tau,v,q)=\langle f,v\rangle.
\end{equation}
{Then the mixed weak formulation has a standard form: find $x\in X$ such that}
\begin{equation}\label{eq:hodge-mixed-abstract-problem}
B(x,y)=F(y),\qquad y\in X.
\end{equation}

% block: hodge-necas-inf-sup page:52 type:proof
{Nečas's generalization of the Lax–Milgram theorem [\BookNumber{61}, Theorem 3.1] tells us when such a problem is well-posed. Well-posedness (for arbitrary $F\in X'$) is equivalent to the \emph{inf-sup condition}}
\begin{equation}\label{eq:hodge-inf-sup-condition}
\gamma:=\inf_{0\ne w\in X}\sup_{0\ne y\in X}\frac{B(w,y)}{\|w\|_X\|y\|_X}>0,
\end{equation}
{in which case the solution $x$ satisfies $\|x\|_X\leq\gamma^{-1}\|F\|_{X'}$. (A second inf-sup condition with the roles of $w$ and $y$ reversed is generally required but in our case is unnecessary, since $B$ is symmetric.) The inf-sup condition is proven in the following theorem.}
\end{Proof}

% block: hodge-inf-sup-theorem page:52 type:theorem
\Needspace{6\baselineskip}
\begin{Theorem}\label{thm:hodge-mixed-inf-sup}
{The bilinear form $B$ defined by \eqref{eq:hodge-mixed-bilinear-form} satisfies the inf-sup condition \eqref{eq:hodge-inf-sup-condition} with $\gamma$ bounded below by a constant depending only on the Poincaré constant $c_P$.}
\end{Theorem}

% block: hodge-inf-sup-proof-task page:52 type:proof
\begin{Proof}
{For any given $(\sigma,u,p)\in V^{k-1}\times V^k\times\mathsf H$ we must exhibit $(\tau,v,q)\in V^{k-1}\times V^k\times\mathsf H$, not identically zero, such that}
\[
B(\sigma,u,p;\tau,v,q)\geq\gamma(\|\sigma\|_V+\|u\|_V+\|p\|)(\|\tau\|_V+\|v\|_V+\|q\|),
\]
{where $\gamma$ depends only on $c_P$. For this, we use the Hodge decomposition $u=u_{\mathcal B}+u_{\mathsf H}+u_{\mathcal B^*}$, where $u_{\mathcal B}=P_{\mathcal B}u$, $u_{\mathsf H}=P_{\mathsf H}u$, and $u_{\mathcal B^*}=P_{\mathcal B^*}u$. Since $u_{\mathcal B}\in\mathcal B^k$, we may write $u_{\mathcal B}=d\rho$ with $\rho\in\mathcal Z^{k-1\perp}$. By the Poincaré inequality,}
\begin{equation}\label{eq:hodge-inf-sup-lifting-bound}
\|\rho\|_V\leq c_P\|u_{\mathcal B}\|.
\end{equation}

% block: hodge-inf-sup-test-choice page:52 type:proof
{Then we define}
\begin{equation}\label{eq:hodge-inf-sup-test-choice}
\begin{gathered}
\tau=\sigma-\frac{1}{c_P^2}\rho\in V^{k-1},\qquad v=-u-d\sigma-p\in V^k,\\
q=-p+u_{\mathsf H}\in\mathsf H^k.
\end{gathered}
\end{equation}

% block: hodge-inf-sup-test-bound page:53 type:proof
{From \eqref{eq:hodge-inf-sup-lifting-bound} and the orthogonality of the Hodge decomposition, we have}
\begin{equation}\label{eq:hodge-inf-sup-test-bound}
\|\tau\|_V+\|v\|_V+\|q\|\leq C(\|\sigma\|_V+\|u\|_V+\|p\|).
\end{equation}
{Plugging \eqref{eq:hodge-inf-sup-test-choice} into \eqref{eq:hodge-mixed-bilinear-form}, we have}
\[
\begin{aligned}
B(\sigma,u,p;\tau,v,q)={}&\|\sigma\|^2+\|d\sigma\|^2+\|du\|^2+\|p\|^2+\|u_{\mathsf H}\|^2\\
&+\frac{1}{c_P^2}\|u_{\mathcal B}\|^2-\frac{1}{c_P^2}\langle\sigma,\rho\rangle.
\end{aligned}
\]

% block: hodge-inf-sup-cross-term page:53 type:proof
{We use the Cauchy–Schwarz inequality, the algebraic-geometric mean inequality, and \eqref{eq:hodge-inf-sup-lifting-bound} to bound the final term:}
\[
\begin{aligned}
\langle\sigma,\rho\rangle&\leq\|\sigma\|\|\rho\|\leq\frac{c_P^2}{2}\|\sigma\|^2+\frac{1}{2c_P^2}\|\rho\|^2\\
&\leq\frac{c_P^2}{2}\|\sigma\|^2+\frac12\|u_{\mathcal B}\|^2.
\end{aligned}
\]
{Thus}
\[
\begin{aligned}
B(\sigma,u,p;\tau,v,q)\geq{}&\frac12\|\sigma\|^2+\|d\sigma\|^2+\|du\|^2+\|p\|^2\\
&+\|u_{\mathsf H}\|^2+\frac{1}{2c_P^2}\|u_{\mathcal B}\|^2.
\end{aligned}
\]

% block: hodge-inf-sup-dual-boundary-bound page:53 type:proof
{Now we apply the Poincaré inequality to $u_{\mathcal B^*}$, for which we have $du_{\mathcal B^*}=du$:}
\[
\|u_{\mathcal B^*}\|_V\leq c_P\|du\|.
\]
{We obtain}
\[
\begin{aligned}
&B(\sigma,u,p;\tau,v,q)\\
&\geq\frac12\|\sigma\|^2+\|d\sigma\|^2+\frac12\|du\|^2+\|p\|^2+\|u_{\mathsf H}\|^2\\
&\qquad+\frac{1}{2c_P^2}\|u_{\mathcal B}\|^2+\frac{1}{2c_P^2}\|u_{\mathcal B^*}\|^2\\
&\geq\eta(\|\sigma\|^2+\|d\sigma\|^2+\|u_{\mathcal B}\|^2+\|u_{\mathsf H}\|^2+\|u_{\mathcal B^*}\|^2+\|du\|^2+\|p\|^2)\\
&=\eta(\|\sigma\|_V^2+\|u\|_V^2+\|p\|^2),
\end{aligned}
\]
{where $\eta=\min(1/2,1/2c_P^2)$ depends only on $c_P$. The inf-sup condition easily follows from this bound and \eqref{eq:hodge-inf-sup-test-bound}.}
\end{Proof}

% block: hodge-general-load-remark page:53 type:remark
\Needspace{6\baselineskip}
\begin{Remark}\label{rem:hodge-general-mixed-load}
{The inf-sup condition in Theorem~\ref{thm:hodge-mixed-inf-sup} implies the well-posedness of the mixed problem \eqref{eq:hodge-mixed-abstract-problem} for general $F\in X'$, not only for $F$ of the form \eqref{eq:hodge-mixed-load-functional}.}
\end{Remark}

% block: hodge-solution-operator page:53 type:ordinary-paragraph
{Having established well-posedness of the Hodge Laplacian, we denote by $K$ the solution operator. That is, for $f\in W$, we define $Kf\in D(L)$ by the equations $LKf=f-P_{\mathsf H}f$, $Kf\perp\mathsf H$. Written in terms of $K$, the solution to the mixed formulation is given by $\sigma=d^*Kf$, $u=Kf$, $p=P_{\mathsf H}f$. Thus the Hodge projections are $P_{\mathcal B}=dd^*K$, $P_{\mathcal B^*}=d^*dK$, and}
\[
u=dd^*Ku+P_{\mathsf H}u+d^*dKu,\qquad u\in W.
\]

% block: hodge-mixed-choice-remark page:53 type:remark
\Needspace{8\baselineskip}
\begin{Remark}\label{rem:hodge-mixed-discretization-motivation}
{The reader may wonder why we chose to establish well-posedness of the abstract Hodge Laplacian using the mixed formulation, rather than the primal weak formulation. After all, the bilinear form for the primal formulation is coercive (by the Poincaré inequality) and hence we can apply the Riesz representation theorem and dispense with the verification of the inf-sup condition. The reason is that we are anticipating the discretization of the Hodge Laplacian, which will be the subject of the next chapter. The discretization will be based on the mixed formulation, which, as we shall see, turns out to be far more appropriate for discretization than the primal formulation. The inf-sup condition for the mixed bilinear form (and not only well-posedness for $F$ of the form \eqref{eq:hodge-mixed-load-functional}) will then be the key to the analysis of the discretization method.}
\end{Remark}

% block: boundary-problems-heading page:54 type:subsection
\subsubsection{{The $\mathcal B$ and $\mathcal B^*$ problems}}\label{subsec:boundary-coboundary-problems}

% block: boundary-problem-hodge-reduction page:54 type:ordinary-paragraph
{In this section, we examine the Hodge Laplace equation $Lv=g$ when the data $g$ belongs to $\mathcal B$ or $\mathcal B^*$. In the former case, $g=df$ for some $f\in V$ (if $g\in V^k$, then $f\in V^{k-1}$). From the solution $u$ of the Hodge Laplace equation}
\begin{equation}\label{eq:hodge-boundary-problem-potential}
Lu=f-P_{\mathsf H}f
\end{equation}
{with data $f$ we obtain the solution to $Lv=g$ by setting $v=du$. This is reasonable since, formally,}
\[
Lv=(dd^*+d^*d)du=dd^*du=d(dd^*+d^*d)u=d(f-P_{\mathsf H}f)=g.
\]

% block: boundary-problem-domain-justification page:54 type:ordinary-paragraph
{This formal calculation is justified as long as $v=du\in D(L)$. To see this, note that, since $u\in D(L)$, $v=du\in V^*$ and $v\in V$. Moreover $dv=0$ and $d^*v=d^*du=P_{\mathcal B^*}f=f-P_{\mathcal Z}f\in V$, since $f\in V$ and $\mathcal Z\subset V$. This shows that $v\in D(L)$, as claimed.}

% block: boundary-problem-intro page:54 type:ordinary-paragraph
{Thus, when the data $g$ for the Hodge Laplacian belongs to $\mathcal B$, the corresponding solution $v$ also belongs to $\mathcal B$. But then the term $d^*dv$ in the Hodge Laplacian vanishes, so $v$ satisfies another problem.}

% block: boundary-problem-formulation page:54 type:ordinary-paragraph
\medskip\noindent\textbf{{The $\mathcal B$ problem}}\par
{Given $g\in\mathcal B^k$, find $v\in\mathcal B^k$ such that $dd^*v=g$.}

% block: boundary-problem-equivalence page:54 type:ordinary-paragraph
{By definition, a solution to this problem must belong as well to $D(dd^*)=\{\,v\in V_k^*\mid d^*v\in V^{k-1}\,\}$. We have just seen that, for $g\in\mathcal B^k$, if $v$ solves the Hodge Laplace problem with data $g$, then $v$ solves the $\mathcal B$ problem. Conversely, if $v$ solves the $\mathcal B$ problem, then we have $dv=0$, and so $Lv=dd^*v=g$, i.e., $v$ solves the Hodge Laplace problem. From the well-posedness of the Hodge Laplace problem, we obtain well-posedness of the $\mathcal B$ problem (but only for $g\in\mathcal B^k$; otherwise there is no solution).}

% block: boundary-problem-well-posedness page:54 type:theorem
\Needspace{6\baselineskip}
\begin{Theorem}\label{thm:hodge-boundary-problem}
{If $g\in\mathcal B^k$, then there is a unique solution $v$ to the $\mathcal B$ problem with data $g$ and this solution coincides with the unique solution to the Hodge Laplace problem with the same data.}
\end{Theorem}

% block: dual-boundary-problem-intro page:54 type:ordinary-paragraph
{Of course, we have the analogous result for $d^*$ as well: if $f\in V^*$ and $u$ the corresponding solution to the Hodge Laplace problem, then $v=d^*u\in D(L)$ and is the solution to the Hodge Laplace problem with data $g=d^*f$. Moreover, if $g\in\mathcal B^*$, then it is also the unique solution to the following problem.}

% block: dual-boundary-problem-formulation page:54 type:ordinary-paragraph
\medskip\noindent\textbf{{The $\mathcal B^*$ problem}}\par
{Given $g\in\mathcal B_k^*$, find $v\in\mathcal B_k^*$ such that $d^*dv=g$.}
